\section{The counterexample}
\label{sec:assembly}

The Hilbert construction supplies the geometric hypotheses of localized
linearization. The resulting completely continuous operator then meets
the weighted graph criterion.

\begin{proof}[Proof of \Cref{thm:main}]
Take the real Hilbert spaces \(U\), \(V=\ell_2\), \(M=U\oplus_2V\)
and the map \(h:M\to U\) from \Cref{thm:hilbert-map}. The blocks
\(U_n\) are finite dimensional, so \(U\) and \(M\) are separable;
as Hilbert spaces they are complete. Put
\[
 E=\F(M),\qquad D=h-\pi_U,\qquad q=\widehat D:E\to U.
\]
The map \(D\) is Lipschitz, \(D(0)=0\), and it has the finite-family
tail-coordinate orthogonality required by
\Cref{thm:localized-linearization}. That theorem makes \(q\)
completely continuous. By \eqref{eq:linearization} and
\eqref{eq:h-bounds},
\[
 \|q\|=\Lip(D)\le\Lip(h)+\|\pi_U\|
 \le\frac{26}{25}+1=\frac{51}{25}.
\]
Set
\[
 Q_1:E\oplus_1\ell_2\to M,\qquad Q_1(e,v)=(qe,v),\qquad
 X=Z_{Q_1},\qquad Y=Z_q.
\]
All hypotheses of \Cref{prop:criterion} hold with
\(a=24/25\), \(b=26/25\). It supplies the coordinatewise bijection
\(\Psi:X\to Y\), an isometric copy of \(c_0(\ell_2)\) in \(X\),
and exclusion of every linearly isomorphic copy from \(Y\).
In particular \(X\) and \(Y\) are not linearly isomorphic.

Since \(L_q\le51/25\) and \(K_q=\max\{L_q,1\}\le51/25\),
\eqref{eq:criterion-bounds} gives
\[
\begin{aligned}
 \Lip(\Psi)&\le
 \max\left\{1+\frac{51}{25},\frac{26}{25}\right\}
 =\frac{76}{25},\\
 \Lip(\Psi^{-1})&\le
 \max\left\{1+\frac{2(51/25)}{24/25},\frac{25}{24}\right\}
 =\frac{21}{4}.
\end{aligned}
\]
The latter inequality is equivalent to the lower bound
\(\|\Psi(s)-\Psi(t)\|_Y\ge(4/21)\|s-t\|_X\).
Both bounds hold for all pairs \(s,t\in X\), exactly as in
\eqref{eq:main-bounds}.
\end{proof}
